\section{Corpus fragments and vocabulary}
Let $\mathcal C$ be a finite corpus of sentences over a finite vocabulary $\Sigma$. We denote $\Sigma^+$ to be the set of finite non-empty strings from vocabulary $\Sigma$. In order to keep the
semantics finite and directly tied to parsing data, let $\Frag(\mathcal C)$ be the set of all non-empty
contiguous strings that occur in at least one sentence in $\mathcal C$. Thus:
\[
\Frag(\mathcal C) \subseteq \Sigma^+.
\]
Because $\mathcal C$ is finite, the set $\Frag(\mathcal C)$ is finite as well.

The choice of $\Frag(\mathcal C)$ rather than all of $\Sigma^+$ is deliberate. It reflects the fact that
this semantics is trained and evaluated relative to observed data. It also ensures that the minima
and maxima appearing in the recursive clauses below are taken over finite sets. Extending the
semantics to arbitrary strings over $\Sigma$ would require a separate treatment of unseen data,
which is left to future work.

